\begin{thebibliography}{99}
\footnotesize\itemsep=0pt

\bibitem{Andrews1967}
Andrews G.E., Enumerative proofs of certain {$q$}-identities, \href{https://doi.org/10.1017/S0017089500000057}{\textit{Glasgow
 Math.~J.}} \textbf{8} (1967), 33--40.

\bibitem{Andrews74}
Andrews G.E., Applications of basic hypergeometric functions, \href{https://doi.org/10.1137/1016081}{\textit{SIAM
 Rev.}} \textbf{16} (1974), 441--484.

\bibitem{Andrews86}
Andrews G.E., The fifth and seventh order mock theta functions, \href{https://doi.org/10.2307/2000275}{\textit{Trans.
 Amer. Math. Soc.}} \textbf{293} (1986), 113--134.

\bibitem{Andrews2012}
Andrews G.E., {$q$}-orthogonal polynomials, {R}ogers--{R}amanujan identities,
 and mock theta functions, \href{https://doi.org/10.1134/S0081543812010038}{\textit{Proc. Steklov Inst. Math.}} \textbf{276}
 (2012), 21--32.

\bibitem{AAR1999}
Andrews G.E., Askey R., Roy R., Special functions, \textit{Encyclopedia Math.
 Appl.}, Vol.~71, \href{https://doi.org/10.1017/CBO9781107325937}{Cambridge University Press}, Cambridge, 1999.

\bibitem{AndrewsBerndt2005}
Andrews G.E., Berndt B.C., Ramanujan's lost notebook. {P}art~{I}, \href{https://doi.org/10.1007/0-387-28124-X}{Springer}, New
 York, 2005.

\bibitem{ADH1988}
Andrews G.E., Dyson F.J., Hickerson D., Partitions and indefinite quadratic
 forms, \href{https://doi.org/10.1007/BF01388778}{\textit{Invent. Math.}} \textbf{91} (1988), 391--407.

\bibitem{AslanIsmail}
Aslan H., Ismail M.E.H., A~{$q$}-translation approach to {L}iu's calculus,
 \href{https://doi.org/10.1007/s00026-019-00450-x}{\textit{Ann. Comb.}} \textbf{23} (2019), 465--488.

\bibitem{ChenWang2020}
Chen D., Wang L., Representations of mock theta functions, \href{https://doi.org/10.1016/j.aim.2020.107037}{\textit{Adv. Math.}}
 \textbf{365} (2020), 107037, 72~pages, \href{https://arxiv.org/abs/1811.07686}{arXiv:1811.07686}.

\bibitem{Gas+Rah}
Gasper G., Rahman M., Basic hypergeometric series, 2nd ed., \textit{Encyclopedia Math.
 Appl.}, Vol.~96, \href{https://doi.org/10.1017/CBO9780511526251}{Cambridge University Press}, Cambridge, 2004.

\bibitem{Jackson1908}
Jackson F., On $q$-functions and a~certain difference operator, \href{https://doi.org/10.1017/S0080456800002751}{\textit{Trans.
 Roy. Soc. Edin.}} \textbf{46} (1909), 253--281.

\bibitem{Krammer1993}
Krammer D., Sums of three squares and {$q$}-series, \href{https://doi.org/10.1006/jnth.1993.1050}{\textit{J.~Number Theory}}
 \textbf{44} (1993), 244--254.

\bibitem{Liu2002}
Liu Z.-G., An expansion formula for {$q$}-series and applications,
 \href{https://doi.org/10.1023/A:1021306016666}{\textit{Ramanujan~J.}} \textbf{6} (2002), 429--447.

\bibitem{Liu2003}
Liu Z.-G., Some operator identities and {$q$}-series transformation formulas,
 \href{https://doi.org/10.1016/S0012-365X(02)00626-X}{\textit{Discrete Math.}} \textbf{265} (2003), 119--139.

\bibitem{Liu2010pacific}
Liu Z.-G., An extension of the quintuple product identity and its applications,
 \href{https://doi.org/10.2140/pjm.2010.246.345}{\textit{Pacific~J.~Math.}} \textbf{246} (2010), 345--390.

\bibitem{Liu2010}
Liu Z.-G., Two {$q$}-difference equations and {$q$}-operator identities,
 \href{https://doi.org/10.1080/10236190902810385}{\textit{J.~Difference Equ. Appl.}} \textbf{16} (2010), 1293--1307.

\bibitem{Liu2011}
Liu Z.-G., An extension of the non-terminating {${}_6\phi_5$} summation and the
 {A}skey--{W}ilson polynomials, \href{https://doi.org/10.1080/10236190903530735}{\textit{J.~Difference Equ. Appl.}} \textbf{17}
 (2011), 1401--1411.

\bibitem{Liu2013RamanJ}
Liu Z.-G., A~{$q$}-series expansion formula and the {A}skey--{W}ilson
 polynomials, \href{https://doi.org/10.1007/s11139-012-9450-2}{\textit{Ramanujan~J.}} \textbf{30} (2013), 193--210.

\bibitem{Liu2013IJTN}
Liu Z.-G., On the {$q$}-derivative and {$q$}-series expansions,
 \href{https://doi.org/10.1142/S1793042113500759}{\textit{Int.~J.~Number Theory}} \textbf{9} (2013), 2069--2089,
 \href{https://arxiv.org/abs/1805.04618}{arXiv:1805.04618}.

\bibitem{LiuRam2013}
Liu Z.-G., On the {$q$}-partial differential equations and {$q$}-series, in The
 Legacy of {S}rinivasa {R}amanujan, \textit{Ramanujan Math. Soc. Lect. Notes
 Ser.}, Vol.~20, Ramanujan Mathematical Society, Mysore, 2013, 213--250,
 \href{https://arxiv.org/abs/1805.02132}{arXiv:1805.02132}.

\bibitem{LiuHahn2015}
Liu Z.-G., A~{$q$}-extension of a partial differential equation and the {H}ahn
 polynomials, \href{https://doi.org/10.1007/s11139-014-9632-1}{\textit{Ramanujan~J.}} \textbf{38} (2015), 481--501,
 \href{https://arxiv.org/abs/1805.07292}{arXiv:1805.07292}.

\bibitem{Liu2017}
Liu Z.-G., On a~system of {$q$}-partial differential equations with applications
 to {$q$}-series, in Analytic Number Theory, Modular Forms and
 {$q$}-hypergeometric Series, \textit{Springer Proc. Math. Stat.}, Vol. 221,
 \href{https://doi.org/10.1007/978-3-319-68376-8_2}{Springer}, Cham, 2017, 445--461, \href{https://arxiv.org/abs/1709.06784}{arXiv:1709.06784}.

\bibitem{Liu2019}
Liu Z.-G., {A}skey--{W}ilson polynomials and a~double {$q$}-series
 transformation formula with twelve parameters, \href{https://doi.org/10.1090/proc/14411}{\textit{Proc. Amer. Math.
 Soc.}} \textbf{147} (2019), 2349--2363, \href{https://arxiv.org/abs/1810.02918}{arXiv:1810.02918}.

\bibitem{Liu2020HRjournal}
Liu Z.-G., A~universal identity for theta functions of degree eight and
 applications, \href{https://doi.org/10.46298/hrj.2021.7432}{\textit{Hardy-Ramanujan~J.}} \textbf{43} (2020), 129--172,
 \href{https://arxiv.org/abs/2104.14705}{arXiv:2104.14705}.

\bibitem{Liu2021}
Liu Z.-G., On the {A}skey--{W}ilson polynomials and a~{$q$}-beta integral,
 \href{https://doi.org/10.1090/proc/15584}{\textit{Proc. Amer. Math. Soc.}} \textbf{149} (2021), 4639--4648.

\bibitem{Liu2023ActaSci}
Liu Z.-G., A~multiple {$q$}-exponential differential operational identity,
 \href{https://doi.org/10.1007/s10473-023-0608-3}{\textit{Acta Math. Sci. Ser.~B (Engl. Ed.)}} \textbf{43} (2023), 2449--2470.

\bibitem{Liu2023Acta}
Liu Z.-G., A~multiple {$q$}-translation formula and its implications,
 \href{https://doi.org/10.1007/s10114-023-2237-0}{\textit{Acta Math. Sin. (Engl. Ser.)}} \textbf{39} (2023), 2338--2363.

\bibitem{Schendel}
Schendel L., Zur {T}heorie der {F}unctionen, \href{https://doi.org/10.1515/crelle-1878-18788406}{\textit{J.~Reine Angew. Math.}}
 \textbf{84} (1878), 80--84.

\bibitem{Shanks1951}
Shanks D., A~short proof of an identity of {E}uler, \href{https://doi.org/10.2307/2032076}{\textit{Proc. Amer. Math.
 Soc.}} \textbf{2} (1951), 747--749.

\bibitem{Shanks1958}
Shanks D., Two theorems of {G}auss, \href{https://doi.org/10.2140/pjm.1958.8.609}{\textit{Pacific~J.~Math.}} \textbf{8}
 (1958), 609--612.

\bibitem{WangChern2019}
Wang C., Chern S., Some {$q$}-transformation formulas and {H}ecke type
 identities, \href{https://doi.org/10.1142/S1793042119500751}{\textit{Int.~J.~Number Theory}} \textbf{15} (2019), 1349--1367.

\bibitem{WangYee2019}
Wang L., Yee A.J., Some {H}ecke--{R}ogers type identities, \href{https://doi.org/10.1016/j.aim.2019.04.015}{\textit{Adv. Math.}}
 \textbf{349} (2019), 733--748.

\end{thebibliography}
